% titel: Blatt 6, A8: Molekülgeometrie nach VSEPR

\subsection*{Lösung}
\begin{center}
\begin{tabular}{@{}l c l@{}}
  Verbindung & EP-Gebiete (Zentralatom) & Molekülgeometrie \\
  \hline
  $\mathrm{O_3}$    & 3 (1 frei)          & gewinkelt (bent), $\approx 117^\circ$ \\
  $\mathrm{AsF_4^-}$& 5 (1 frei)          & see-saw (Wippe) \\
  $\mathrm{BrF_5}$  & 6 (1 frei)          & quadratisch-pyramidal \\
  $\mathrm{ICl_4^-}$& 6 (2 frei)          & quadratisch-planar \\
  $\mathrm{SnH_4}$  & 4 (0 frei)          & tetraedrisch \\
  $\mathrm{XeF_3^+}$& 5 (2 frei)          & T-förmig \\
  $\mathrm{AsH_3}$  & 4 (1 frei)          & trigonal-pyramidal \\
  $\mathrm{IBr_2^-}$& 5 (3 frei)          & linear \\
\end{tabular}
\end{center}

\subsection*{Kurzbegründungen}
\begin{itemize}
  \item $\mathrm{O_3}$: zentrales O mit einer Doppel-, einer Einfachbindung und
        einem freien Paar $\to$ 3 Gebiete, gewinkelt.
  \item $\mathrm{AsF_4^-}$ / $\mathrm{XeF_3^+}$: As hat 5 Gebiete (4 Bindungen +
        1 frei) $\to$ see-saw; $\mathrm{XeF_3^+}$ hat wegen der $+$-Ladung ein
        Valenzelektron weniger, also 3 Bindungen + 2 freie Paare $\to$ T-förmig.
  \item $\mathrm{BrF_5}$ (6 Gebiete, 1 frei) $\to$ quadratisch-pyramidal;
        $\mathrm{ICl_4^-}$ (6 Gebiete, 2 frei) $\to$ quadratisch-planar.
  \item $\mathrm{SnH_4}$ (4 Gebiete, keine frei) $\to$ tetraedrisch;
        $\mathrm{AsH_3}$ (4 Gebiete, 1 frei) $\to$ trigonal-pyramidal;
        $\mathrm{IBr_2^-}$ (5 Gebiete, 3 frei) $\to$ linear.
\end{itemize}

\begin{ergebnis}
geometrien: $\mathrm{O_3}$ bent, $\mathrm{AsF_4^-}$ see-saw, $\mathrm{BrF_5}$
quadratisch-pyramidal, $\mathrm{ICl_4^-}$ quadratisch-planar, $\mathrm{SnH_4}$
tetraedrisch, $\mathrm{XeF_3^+}$ T-förmig, $\mathrm{AsH_3}$ trigonal-pyramidal,
$\mathrm{IBr_2^-}$ linear.
\end{ergebnis}

\begin{kontrolle}
Die EP-Gebiete (3,5,6,6,4,5,4,5) und die freie-Paare-Zahl bestimmen eindeutig
jede Zeile; spaltenweise Prüfung der Summe (Bindungen + freie = Gebiete) füllt
die Tabelle. Das deckt sich mit den bekannten beobachteten Formen dieser
Verbindungen.
\end{kontrolle}

\begin{merksatz}
Zuerst EP-Gebiete zählen, dann die freien Paare ``abrechnen''. 5 Gebiete:
see-saw (1 frei), T-förmig (2), linear (3); 6 Gebiete: pyramidal (1),
quadratisch-planar (2).
\end{merksatz}
